\documentclass[11pt]{article}

\usepackage[margin=1.1in]{geometry}
\usepackage{amsmath,amssymb,amsthm}
\usepackage{booktabs}
\usepackage{graphicx}
\usepackage{microtype}
\usepackage[table]{xcolor}
\usepackage{caption}
\usepackage{enumitem}
\usepackage[hidelinks]{hyperref}
\usepackage{natbib}


\newtheorem{proposition}{Proposition}
\newtheorem{corollary}[proposition]{Corollary}

\newcommand{\weff}{w_{\mathrm{eff}}}
\newcommand{\dgeo}{d_{\mathrm{geo}}}
\newcommand{\hold}{\textsf{hold}}
\newcommand{\swi}{\textsf{switch}}

\title{\bfseries The Mechanism Is Not the Eponym:\\
Budget-Matched Ablation of Four Classical AI Method Families}

\author{%
  Muhaiminul Islam Ninad\\
  \small Department of Computer Science and Engineering\\
  \small University of Dhaka\\
  \small \texttt{muhaiminulislam-2021111219@cs.du.ac.bd}
}

\date{\today}

\begin{document}
\maketitle

\begin{abstract}
Classical AI method families are named for a mechanism --- the \emph{heuristic} in heuristic
search, the \emph{constraints} in constraint propagation, the \emph{swarm} in particle swarm
optimisation, the \emph{learning} in reinforcement learning --- and both textbooks and practice
credit that mechanism with the family's performance. We ask, for four families at once, whether
the eponymous mechanism is the mechanism that does the work. We build four problems in a single
application setting (urban mobility in Dhaka, Bangladesh), implement each family from scratch, and
run ablations under strictly matched compute budgets, scoring each family on both its own
customary metric and on a metric that charges every resource consumed.

In no family does the standard attribution survive the ablation intact, and in three of them the
customary metric is what conceals the failure. \textbf{(i)~Informed search:} the three ``different'' heuristics in a standard
routing study are positive scalar multiples of one function, so greedy best-first search is
provably invariant to the choice (confirmed: identical paths on 10/10 queries) and $\mathrm{A}^*$
over them is a one-parameter weight sweep whose node counts and search times are perfectly rank-correlated
with that scalar ($\rho = -1.0000$). Nodes-expanded ranks $\mathrm{A}^*$ above uniform-cost search
by $2.15\times$; wall clock ranks it \emph{below} by $5.98\times$, because $88.6\%$ of
$\mathrm{A}^*$'s runtime is a one-off $O(|E|)$ heuristic calibration that the metric prices at zero. The fastest exact method in the study is
uninformed. \textbf{(ii)~Constraint satisfaction:} ordering and propagation reduce node counts
exactly as folklore predicts ($-3$ to $-5$ nodes, $p \le 10^{-4}$) while running $12$--$14\times$
slower in wall clock and changing which instances are solvable in 1 case out of 80 --- that one
adversely; past the point where search stops completing, the node metric inverts and measures
throughput rather than pruning. \textbf{(iii)~Population-based search:} thirty non-communicating
particles are statistically indistinguishable from random search ($p = 0.22$) and a single particle
is \emph{worse} than random search ($p < 10^{-4}$); switching on the social term alone, at the same
budget and the same random stream, buys $7.6$ minutes of every student's expected waiting time. The
population is not the mechanism, the communication channel is.
\textbf{(iv)~Reinforcement learning:} on identical data --- same seeds, same 400{,}000 transitions,
same order --- planning on an estimated model beats a \emph{tuned} model-free learner ($1.74\%$
versus $2.72\%$ regret; $89.9\%$ versus $65.7\%$ action-optimality), and planning with arrival rates
wrong by a factor of ten still beats learning from scratch.

We give a common account of the four results: in each family the operative component is one whose
cost is paid once and reused many times --- an amortisable precomputation or a broadcast
information channel --- while the eponym's cost is paid per unit of work, and each subfield's
customary metric is precisely the one that declines to charge for reuse. We close with a
metric-audit rule and release the testbed, code, and every number in this paper.
\end{abstract}

\section{Introduction}

A method family's name is a hypothesis about what makes it work. ``Heuristic search'' asserts that
the heuristic is what turns an intractable search into a tractable one. ``Constraint propagation''
asserts that propagating constraints is what prunes the tree. ``Particle \emph{swarm}
optimisation'' asserts that the swarm --- a population of interacting searchers --- is the source
of the optimisation power. ``Reinforcement \emph{learning}'' asserts that learning from experience
is the right response to not having a model.

These are good hypotheses. They are also rarely tested \emph{as} hypotheses, because the natural
experiment --- remove the named mechanism, hold everything else fixed, and see what happens ---
requires a budget-matching discipline that the field's customary reporting metrics do not enforce
and sometimes actively frustrate. Heuristic search reports nodes expanded, a metric that charges
nothing for computing the heuristic. Metaheuristics report generations or iterations, metrics
denominated in units whose cost varies with population size. Reinforcement learning reports
sample complexity, a metric that charges nothing for what is done with a sample after it is
collected. In each case the customary metric is the one under which the eponymous mechanism looks
best, and this is not a coincidence: metrics are chosen by the communities that the mechanism
defines.

This paper runs the natural experiment for four families at once. We construct four problems in a
single application setting --- urban mobility in Dhaka, a megacity for which almost no operational
data is published --- so that the four ablations share a domain, an implementation discipline, and
an author. Each family is implemented from scratch, without solver libraries, so that the ablations
can reach inside the method. Each comparison is budget-matched: identical objective-evaluation
counts for the population methods, identical transition budgets for the learners, identical
wall-clock budgets for the constraint solvers, and, for search, both the customary metric and a
metric that charges every millisecond spent.

\paragraph{Findings.} In each family the eponym is either not the operative mechanism, or is
operative only under a metric that declines to charge what it costs. We state each result with its
test in Sections~\ref{sec:search}--\ref{sec:rl}. The headline in each case:

\begin{itemize}[leftmargin=1.4em, itemsep=2pt]
\item \textbf{Search.} The heuristic ``portfolio'' in a conventional multi-heuristic routing study
  collapses to a single scalar, and the standard metric misprices the heuristic's cost by more than
  an order of magnitude. The fastest exact method we measure is an uninformed one.
\item \textbf{Constraints.} Ordering and propagation prune as advertised and cost an order of
  magnitude more than they save, and once instances stop being solvable within budget the node
  metric measures the opposite of what it appears to. Repair-based search moves the cost of an
  attempt by three orders of magnitude while abandoning the mechanism the family is named for.
\item \textbf{Population.} Neither the population nor the individual is doing the optimisation. A
  single one-bit-per-iteration broadcast channel is.
\item \textbf{Learning.} The model is doing the work, not the learning; and the model can be
  grossly wrong and still win.
\end{itemize}

\paragraph{Contributions.}
\begin{enumerate}[leftmargin=1.6em, itemsep=2pt]
\item A budget-matched ablation protocol applied uniformly to four classical AI families
  (\S\ref{sec:protocol}), with a from-scratch reference implementation of each.
\item Four ablation results (\S\ref{sec:search}--\S\ref{sec:rl}), three of which disconfirm the
  standard attribution and one of which (\S\ref{sec:rl}) disconfirms it after we first
  \emph{strengthen} the arm we expected to lose --- a step that reduces our own headline effect
  sevenfold and which we report because omitting it would have been the more publishable choice.
\item A cross-family account of \emph{why} the credit is misassigned (\S\ref{sec:synthesis}): in
  every case the operative component is one whose cost is amortisable across units of work, and
  every customary metric is one that does not charge for amortisation. We give a corresponding
  metric-audit rule.
\item A released, deterministic testbed: four problems, four reference implementations, 51 tests,
  and scripts that regenerate every table in this paper.
\end{enumerate}

\paragraph{What this paper is not.} It is not a claim that heuristics, propagation, populations,
or model-free learning are useless; all four are load-bearing under conditions we identify. It is
not a claim about modern learned systems --- every method here is tabular, exact, or small.
And it is not a benchmark competition: no method here is proposed as state of the art.
Section~\ref{sec:threats} is unusually long by design.

\section{Related Work}

\paragraph{Ablation as critique.} A recurring genre establishes that a field's headline progress
survives less well than reported once baselines are tuned and budgets matched: neural language
models \citep{melis2018state}, deep reinforcement learning \citep{henderson2018deep}, metric
learning \citep{musgrave2020metric}, and neural recommendation \citep{dacrema2019areweb}. The
methodological diagnoses --- unmatched tuning budgets, weak baselines, metric choice --- are
collected by \citet{lipton2019troubling} and \citet{sculley2018winners}. Our contribution differs
in scope rather than in kind: instead of auditing one family's recent literature, we ablate the
defining mechanism of four families simultaneously on one testbed, which lets us ask whether the
misattributions share a structure. They do (\S\ref{sec:synthesis}).

\paragraph{Experimental algorithmics.} That algorithm comparisons must charge all resources and
report them honestly is old advice \citep{hooker1995testing, johnson2002theoretician,
mcgeoch2012guide}. \citet{hooker1995testing} in particular argues for controlled experiments that
isolate mechanisms rather than competitive testing of implementations, which is precisely the
design here. \citet{burns2012implementing} document how far heuristic-search wall-clock results
depend on implementation constants --- a caution we take seriously in \S\ref{sec:threats}, since
one of our results is a wall-clock reversal.

\paragraph{Search.} $\mathrm{A}^*$ \citep{hart1968formal} and its optimality theory
\citep{dechter1985generalized, pearl1984heuristics} are stated in terms of node expansions;
weighted $\mathrm{A}^*$ \citep{pohl1970heuristic} trades expansions for bounded suboptimality.
Our \S\ref{sec:search} shows a standard multi-heuristic study reducing exactly to
\citeauthor{pohl1970heuristic}'s one-parameter family, and that the expansion-count metric that
underwrites the theory misprices the heuristic in practice. Strong uninformed baselines --- here
bidirectional uniform-cost search, in the tradition of \citet{holte2016bidirectional} --- are what
make the reversal visible.

\paragraph{Constraint satisfaction.} Forward checking \citep{haralick1980increasing} and arc
consistency \citep{mackworth1977consistency} define the propagation family; minimum-remaining-values
and least-constraining-value orderings are its standard companions. Min-conflicts
\citep{minton1992minimizing} is the repair-based alternative, historically presented as the
surprising outsider. \S\ref{sec:csp} finds the outsider dominating on every operational metric.

\paragraph{Population-based search.} PSO \citep{kennedy1995particle} with inertia weighting
\citep{shi1998modified}, its topology literature \citep{kennedy2002population}, and real-coded GAs
with BLX-$\alpha$ \citep{eshelman1993real} are all implemented here from their original update
rules. \citet{sorensen2015metaheuristics} argues that metaphor-driven metaheuristics obscure their
own mechanisms; \S\ref{sec:pso} supplies a quantitative instance, isolating the social term as the
entire source of PSO's advantage on our problem. \citet{wolpert1997no} frames why such results must
be stated relative to a problem class.

\paragraph{Decision making under uncertainty.} Value iteration \citep{bellman1957dynamic,
puterman1994markov} and $Q$-learning \citep{watkins1992qlearning, sutton2018reinforcement} are the
model-based and model-free poles. The intermediate position --- estimate a model from experience,
then plan on the estimate --- runs from Dyna \citep{sutton1990integrated} through the PAC-MDP
analyses \citep{kearns2002near, brafman2002rmax} to certainty-equivalence results in control
\citep{mania2019certainty}. That model-based methods can be more sample-efficient is well
established in theory; \S\ref{sec:rl} contributes a controlled empirical instance in which the
comparison is run on \emph{byte-identical} data against a \emph{tuned} model-free baseline, and in
which the model is then degraded until it loses, locating the crossing point.

\paragraph{Domain.} The road network is from OpenStreetMap via OSMnx \citep{boeing2017osmnx}. The
signal-control baseline is Webster's fixed-time method \citep{webster1958traffic}; delay accounting
uses Little's law \citep{little1961proof}.

\section{Testbed and Protocol}
\label{sec:protocol}

\subsection{Four problems, one setting}

Table~\ref{tab:testbed} summarises the testbed. The four problems were designed independently as
instances of four method families, and share only their setting: mobility in Dhaka, and for the
last two specifically the University of Dhaka campus. Sharing a setting is a convenience for
exposition and a control on author effects; it is not a claim that the four problems are otherwise
related. Problems~3 and~4 in particular are deliberately different in kind --- an offline design
problem over a continuous decision vector, and an online sequential control problem under
uncertainty.

\begin{table}[t]
\centering
\small
\caption{The four problems, the family each instantiates, the mechanism whose credit is under test,
and the resource held fixed across arms.}
\label{tab:testbed}
\begin{tabular}{@{}llll@{}}
\toprule
\textbf{Problem} & \textbf{Family} & \textbf{Eponym under test} & \textbf{Matched budget} \\
\midrule
Multi-factor routing      & Informed search        & the heuristic        & per-query wall clock \\
Fuel-crisis allocation    & Constraint satisfaction& ordering + propagation & 5\,s per instance \\
Shuttle timetabling       & Population-based search& the population       & 3{,}030 objective evaluations \\
Signal control            & Reinforcement learning & the learning         & 400{,}000 transitions \\
\bottomrule
\end{tabular}
\end{table}

\paragraph{Problem 1 --- multi-factor routing.} Shortest-path queries on the Dhaka OpenStreetMap
drive network (55,009 nodes, 137,529 directed edges) under a multi-factor edge cost that
combines length, road class, congestion, surface quality, and a traveller-context term. Six search
algorithms (uniform-cost, Dijkstra, bidirectional uniform-cost, $\mathrm{A}^*$, weighted
$\mathrm{A}^*$, greedy best-first) run over three named heuristics, giving twelve configurations on
each of ten randomly drawn landmark pairs.

\paragraph{Problem 2 --- fuel-crisis allocation.} Assigning $n$ vehicles to
(station, pump, time-slot) triples under fuel-type compatibility, vehicle range, pump exclusivity,
station supply, and time windows, minimising a weighted sum of travel distance, waiting time, a
quadratic lateness penalty for emergency vehicles, and unassigned-vehicle count. Five solvers:
chronological backtracking; backtracking with minimum-remaining-values (MRV) and degree tie-break;
with least-constraining-value (LCV); with forward checking plus MRV; and min-conflicts local search
with a tabu list.

\paragraph{Problem 3 --- shuttle timetabling.} Choosing $K = 14$ continuous departure times in a
780-minute service window to minimise mean student waiting time, given bus capacity $C = 40$,
fleet size $B = 3$, non-homogeneous Poisson arrivals with class-change bursts, congestion-dependent
round-trip times, reneging after 60 minutes, and a quadratic fleet-overload penalty. The objective
is non-differentiable, piecewise constant, and invariant under permutation of the decision vector
(so every optimum has $14! \approx 8.7\times 10^{10}$ symmetric copies). Solvers: PSO and a
real-coded GA, both from scratch, plus four baselines.

\paragraph{Problem 4 --- signal control.} A finite discounted continuing MDP for a two-approach
signalised intersection: state $(q_A, q_B, \phi, e, p)$ with queues truncated at 12, phase
indicator, phase timer, and an exogenous traffic regime, giving $|\mathcal{S}| = 7,098$ and
11,154 legal state--action pairs under minimum- and maximum-green masks. Actions are $\hold$
and $\swi$; switching burns a clearance tick; reward charges queue length, switching cost, and
spillback. Solvers: value iteration, $Q$-learning, and certainty-equivalence planning, all from
scratch, plus five baselines including Webster fixed-time.

\subsection{Protocol}

\paragraph{Budget matching is the load-bearing control.} Within each family, every arm receives an
identical allocation of the family's scarce resource (Table~\ref{tab:testbed}, last column). In
Problem~3 this is enforced structurally: the objective is callable from exactly one function per
solver file, that function increments a counter, and the counter is asserted equal to
$30 \times (1 + 100) = 3030$ for every arm including every ablation. In Problem~4, the
certainty-equivalence planner is handed the \emph{literal transition list} that $Q$-learning
generated --- same seed, same transitions, same order --- so that any difference is attributable to
what each method does with the data rather than to how much it received.

\paragraph{Determinism.} Every experiment takes an explicit seed; there is no wall-clock or
unseeded randomness in any solver. Re-running any experiment reproduces its table bit-for-bit.
Where an arm is stochastic we report paired results over seeds and test them.

\paragraph{Statistics.} Comparisons over seeds use the Wilcoxon signed-rank test
\citep{wilcoxon1945individual} for location and Levene's test \citep{levene1960robust} for
dispersion, on 30 paired seeds for Problem~3 and 5 for Problem~4. We report $p$-values and make no
claim of ``better'' without a test. We report disconfirmations of our own expectations in the same
register as confirmations; three of the four sections below contain at least one.

\paragraph{Provenance.} The four problems and their reference implementations originate as
independent coursework assignments, each specified before its results were known and each carrying
a written instruction to report disconfirming results rather than retune the instance. That origin
is worth stating plainly: it explains the from-scratch implementations (no solver libraries were
permitted), the per-problem test suites, and the fact that the four families were not selected
after the fact to produce a common conclusion --- the ablations in \S\ref{sec:pso} and
\S\ref{sec:rl} were specified as experiments before either was run. It also bounds the scale:
these are course-sized instances, not industrial benchmarks (\S\ref{sec:threats}). The
cross-family reading in \S\ref{sec:synthesis} is, by contrast, entirely post hoc.

\paragraph{Reporting only what runs.} Every number in this paper was regenerated from the released
code by the scripts in \texttt{repro/}, on the hardware in Appendix~\ref{app:env}. Where a method
is described in project documentation but is not present and executable in the released code, it
is excluded from this paper entirely rather than reported from a stale artifact.

\section{Informed Search: The Heuristic Is One Scalar, and Nodes Expanded Does Not Charge For It}
\label{sec:search}

\subsection{The heuristic portfolio collapses to one parameter}

The study defines three named heuristics --- \texttt{admissible}, \texttt{fast}, and
\texttt{realism} --- and crosses them with three informed algorithms, in the standard pedagogical
pattern of ``one admissible heuristic and two inadmissible ones.'' All three have the form
\begin{equation}
h_i(n) \;=\; k_i \cdot \dgeo(n, \text{goal}), \qquad k_i > 0,
\label{eq:hform}
\end{equation}
where $\dgeo$ is great-circle distance and $k_i$ is a constant depending on the graph, the
traveller context, and the cost preset but \emph{not} on $n$. Concretely,
$k_{\text{adm}} = m^{*}$, the global minimum cost per metre;
$k_{\text{fast}} = 1.25\, m^{*}$; and $k_{\text{realism}} = 1.08\,\bar{m}$ where $\bar{m}$ is the
\emph{mean} cost per metre. Measured on this graph, $m^{*} = 1.0113$ and $\bar{m} = 2.3213$, so
$\bar{m}/m^{*} = 2.295$.

Two consequences follow immediately.

\begin{proposition}[Greedy best-first is scale-invariant]
\label{prop:greedy}
Greedy best-first search orders its frontier by $h$ alone. For $k > 0$, the orderings induced by
$h$ and $k\,h$ are identical. Hence greedy best-first with any two heuristics of the form
\eqref{eq:hform} expands the same nodes in the same order and returns the same path.
\end{proposition}

\begin{proposition}[The informed configurations are one weight sweep]
\label{prop:weight}
$\mathrm{A}^*$ with $f = g + k_i\,\dgeo$ is exactly weighted $\mathrm{A}^*$
\citep{pohl1970heuristic} with $f = g + \weff \cdot h_{\text{adm}}$ at weight
$\weff = k_i / m^{*}$; and weighted $\mathrm{A}^*$ at explicit weight $w$ over $h_i$ is the same
algorithm at $\weff = w\,k_i/m^{*}$. The nine informed configurations therefore occupy nine points
on a single scalar axis, with uniform-cost search at $\weff = 0$ and greedy best-first at
$\weff = \infty$.
\end{proposition}

Proposition~\ref{prop:greedy} is confirmed exactly: across all ten queries, greedy best-first
returns \emph{byte-identical paths} under all three heuristics (10/10), expanding 120 nodes and
incurring $30.754\%$ suboptimality in every case. The three-heuristic portfolio is, for this
algorithm, one heuristic.

Table~\ref{tab:weff} orders every configuration by $\weff$. Both node count and search time are
\emph{perfectly} rank-correlated with the scalar (Spearman $\rho = -1.0000$ for each over the six
distinct finite-weight configurations), and suboptimality is monotone non-decreasing
($\rho = +0.845$, with four exact ties at zero). There is no second dimension in this design: what
was presented as a $3 \times 3$ grid of algorithms and heuristics is a one-dimensional sweep of
$\weff$, and everything measured is a monotone function of it.

\begin{table}[t]
\centering
\small
\caption{All twelve configurations ordered by effective weight $\weff$ (Prop.~\ref{prop:weight}).
Means over ten landmark pairs. \emph{Search} excludes heuristic setup; \emph{total} includes it.
Gap is cost above optimal. Node count and search time are perfectly rank-correlated with $\weff$.}
\label{tab:weff}
\setlength{\tabcolsep}{4.5pt}
\begin{tabular}{@{}lrrrrrr@{}}
\toprule
\textbf{Configuration} & $\weff$ & \textbf{Nodes} & \textbf{Search (ms)} & \textbf{Setup (ms)} &
\textbf{Total (ms)} & \textbf{Gap (\%)} \\
\midrule
Uniform-cost search        & 0     & 14,735 & 295.3 & 0.8    & 296.0  & 0.000 \\
Dijkstra                   & 0     & 14,735 & 303.4 & 0.6    & 304.0  & 0.000 \\
Bidirectional UCS          & 0     & 7,316  & 144.1 & 0.4    & \textbf{144.6}  & 0.000 \\
$\mathrm{A}^*$ + admissible          & 1.000 & 6,845  & 202.6 & 1567.7 & 1770.3 & 0.000 \\
$\mathrm{A}^*$ + fast                & 1.250 & 5,518  & 152.8 & 1552.2 & 1705.0 & 0.000 \\
Weighted $\mathrm{A}^*$ + admissible & 1.350 & 5,057  & 117.4 & 1574.4 & 1691.7 & 0.000 \\
Weighted $\mathrm{A}^*$ + fast       & 1.688 & 3,540  & \textbf{81.4}  & 1544.3 & 1625.7 & 0.000 \\
$\mathrm{A}^*$ + realism             & 2.479 & 1,238  & 58.5  & 1555.3 & 1613.8 & 0.080 \\
Weighted $\mathrm{A}^*$ + realism    & 3.347 & 492   & 16.7  & 1563.1 & 1579.7 & 3.222 \\
Greedy best-first + admissible & $\infty$ & 120 & 3.0 & 1567.3 & 1570.3 & 30.754 \\
Greedy best-first + realism    & $\infty$ & 120 & 2.9 & 1516.3 & 1519.2 & 30.754 \\
Greedy best-first + fast       & $\infty$ & 120 & 3.2 & 1520.9 & 1524.1 & 30.754 \\
\bottomrule
\end{tabular}
\end{table}

This is not a defect peculiar to this testbed; it is the default outcome of the standard
instruction to supply ``one admissible and two inadmissible heuristics'' for a geometric domain,
where the only readily available admissible quantity is a scaled straight-line distance and the
natural way to make it inadmissible is to scale it further. The lesson is that a heuristic
portfolio must be checked for functional independence before it is crossed with anything, and that
the check is cheap: compute the induced orderings.

\subsection{The metric reversal}

Node expansions are the customary metric of heuristic search, and by that metric
$\mathrm{A}^*$ with the admissible heuristic dominates uniform-cost search: 6,845 versus
14,735 expansions, a $2.15\times$ reduction. By total wall clock the ranking reverses:
1770.3\,ms versus 296.0\,ms, a $5.98\times$ \emph{loss}.

The reversal has a single cause, which we measure rather than infer. Both heuristic constants
$m^{*}$ and $\bar{m}$ are global reductions over the edge set, computed by two full $O(|E|)$ passes
that evaluate the multi-factor cost function on all 137,529 edges. Timed in isolation over
five repetitions this costs 1,554--2,021\,ms (median 1,725\,ms);
measured inside the batch it accounts for 1,516--1,574\,ms of every informed
query, against 0.4--0.8\,ms for the uninformed methods. It accounts for
$88.6\%$ of $\mathrm{A}^*$'s total runtime, against $0.3\%$ for uniform-cost search, and the node-count
metric charges exactly none of it.

Charging it correctly changes three conclusions.

\paragraph{(i) Even excluding setup, node counts overstate the heuristic's benefit by $1.48\times$.}
$\mathrm{A}^*$ expands $2.15\times$ fewer nodes but spends only $1.46\times$ less time searching
(202.6\,ms versus 295.3\,ms), because each $\mathrm{A}^*$ expansion
costs 29.60\,$\mu$s against uniform-cost's 20.04\,$\mu$s --- the heuristic
is evaluated on every generated node. A node is not a unit of cost; it is a unit of cost only if
the per-node cost is equal across arms, which is exactly what introducing a heuristic violates.

\paragraph{(ii) The advantage is real but amortisable, and the break-even is quantifiable.}
The setup constant depends only on the graph, the traveller context, and the cost preset --- not on
the query. Hoisting it out of the query loop, $\mathrm{A}^*$+admissible saves
92.7\,ms of search per query against uniform-cost search, so the
1567.7\,ms setup pays for itself after $16.9$ queries sharing a context. Against the
stronger uninformed baseline the arithmetic is harsher: bidirectional uniform-cost search searches
in 144.1\,ms, faster than $\mathrm{A}^*$+admissible's 202.6\,ms, so
$\mathrm{A}^*$ with the admissible heuristic \emph{never} breaks even against it at any query
volume. The first informed configuration that does is weighted $\mathrm{A}^*$+fast, at
81.4\,ms of search, which saves 62.7\,ms per query and breaks even
after $24.6$ queries.

\paragraph{(iii) The fastest exact method in the study is uninformed.} Bidirectional uniform-cost
search returns provably optimal paths on 10/10 queries in 144.6\,ms total ---
$11.2\times$ faster than the fastest exact informed configuration and $12.2\times$ faster than
$\mathrm{A}^*$ with the admissible heuristic. It achieves a $2.01\times$ node reduction over
unidirectional uniform-cost search with no precomputation at all: a structural saving rather than
an informational one, and one that the query-volume argument in (ii) does not erode.

\subsection{Verdict}

The heuristic in this study is one scalar, and its benefit is real but is (a) overstated
$1.48\times$ by node counting even before setup, (b) overwhelmed for single queries by an
$O(|E|)$ precomputation that node counting prices at zero, and (c) in its admissible form, beaten
at \emph{any} query volume by an uninformed algorithm exploiting problem structure. Only an
inadmissible configuration --- weighted $\mathrm{A}^*$ at $\weff = 1.688$ --- eventually overtakes
bidirectional search, and then only past 25 queries per context. The mechanism that most improves
single-query exact routing here is bidirectional expansion, not information.

\section{Constraint Satisfaction: The Orderings Do Not Change What Gets Solved}
\label{sec:csp}

Constraint satisfaction credits its performance to two mechanisms that share the family's name:
\emph{ordering} --- choosing which variable to instantiate next (minimum remaining values, degree
tie-break) and which value to try first (least constraining value) --- and \emph{propagation} ---
pruning domains after each assignment (forward checking). Both are taught as the difference between
an exponential blow-up and a tractable search.

We ran five solvers --- chronological backtracking, {\tt bt\_mrv}, {\tt bt\_lcv},
{\tt bt\_fc\_mrv\_deg} (forward checking + MRV + degree), and min-conflicts local search
\citep{minton1992minimizing} --- over 10 instance sizes $\times$ 8 seeds at a 5\,s budget per
instance: 400 runs. Every run records whether it exhausted the budget, which turns out to matter
more than anything else in this section.

\subsection{The orderings do not change what gets solved}

Table~\ref{tab:csp-solved} gives the outcome that the ordering literature would most expect to
move, and it does not move. All four systematic solvers solve the same number of instances to
within one, and pairwise they disagree on the outcome of \textbf{1 of 80 instances} --- at $n = 25$,
seed 303, where {\tt bt\_lcv} fails and the other three succeed. That single disagreement runs
\emph{against} the heuristic. Adding MRV, LCV, or forward checking to chronological backtracking
changed which instances were solvable, at this budget, essentially not at all.

\begin{table}[t]
\centering
\small
\setlength{\tabcolsep}{3.5pt}
\caption{Left: outcomes over all 400 runs. Right: the 24 instances every solver solved, where node
counts are uncensored and therefore meaningful; medians there, means on the left. Paired Wilcoxon
against chronological backtracking, on node counts.}
\label{tab:csp-solved}
\begin{tabular}{@{}lrrrrrl@{}}
\toprule
& \multicolumn{3}{c}{\textbf{all 400 runs}} & \multicolumn{3}{c}{\textbf{24 commonly solved instances}} \\
\cmidrule(lr){2-4}\cmidrule(lr){5-7}
\textbf{Solver} & \textbf{solved} & \textbf{nodes} & \textbf{time} &
\textbf{nodes} & \textbf{time} & \textbf{vs.\ basic} \\
\midrule
Chronological backtracking & 27/80 & 444,111 & 3.15\,s & 13.5 & 0.09\,ms & --- \\
\quad + MRV / degree       & 27/80 & 322,997 & 3.02\,s & 13.5 & 0.14\,ms & $+0.0$, $p = 0.27$ \\
\quad + LCV                & 26/80 & 223,080 & 3.23\,s & 10.5 & 1.09\,ms & $-3.0$, $p = 10^{-4}$ \\
\quad + FC + MRV + deg.\ & 27/80 & 126,481 & 3.17\,s & \phantom{0}9.0 & 1.27\,ms & $-5.0$, $p < 10^{-4}$ \\
\midrule
Min-conflicts (repair) & 25/80 & \textbf{119} & \textbf{0.044\,s} & \textbf{1.0} & \textbf{0.08\,ms} & $-13.5$, $p < 10^{-4}$ \\
\bottomrule
\end{tabular}
\end{table}

\subsection{The orderings do reduce nodes --- and make the solver slower}

On the 24 instances every solver solved --- the only regime where a node count means what it is
supposed to mean --- the ordering heuristics behave exactly as folklore predicts, and the
improvement is statistically solid. Least-constraining-value removes a median of 3.0 nodes
($p = 10^{-4}$) and forward checking with MRV removes 5.0 ($p < 10^{-4}$), against a median of 13.5
for chronological backtracking. MRV on its own is the exception: a median difference of exactly
zero ($p = 0.27$).

Every one of them is slower. LCV takes $12.1\times$ and forward checking $14.1\times$ the wall clock
of chronological backtracking on the same instances (1.09\,ms and 1.27\,ms against 0.09\,ms), and
even MRV --- which saves no nodes at all --- costs $1.6\times$. The per-node price of computing the
ordering and running the propagation exceeds, by an order of magnitude, the value of the nodes it
saves. This is the same structure as \S\ref{sec:search}: the customary metric shows the eponymous
mechanism working, and it \emph{is} working, but the metric does not charge what it costs to run.

The one arm that improves both metrics simultaneously is the one that abandons systematic search.
Min-conflicts expands a median of 1.0 node at 0.08\,ms --- as fast as chronological backtracking
per attempt and $13.5$ nodes cheaper --- and over all 400 runs averages 119 nodes and 0.044\,s
against 444,111 nodes and 3.15\,s: $3{,}700\times$ fewer nodes and $72\times$ less time, for
25 solved instances against 27. Both solve rates are low in absolute terms --- $31\%$ and $34\%$ of
each solver's 80 runs, because most instances in this sweep are out of reach for every solver at a
5\,s budget --- but min-conflicts reaches $93\%$ of chronological backtracking's solve count for $0.03\%$ of its
search cost, by carrying one complete (initially infeasible) assignment forward and repairing it
rather than rebuilding a partial assignment prefix on every backtrack.

\subsection{Under censoring, the node metric inverts}

At $n \geq 30$ every systematic run exhausts the 5\,s budget and nothing is solved. Aggregate node
counts over such runs are widely reported anyway --- our own source artifact reports them --- and
they are worse than uninformative. When every solver stops at the same clock, its node count
measures \emph{how fast it grinds}, not how well it prunes:

\begin{center}
\small
\begin{tabular}{@{}lrrr@{}}
\toprule
\textbf{Solver} & \textbf{$n=30$} & \textbf{$n=40$} & \textbf{$n=50$} \\
\midrule
Chronological backtracking & 163,065 & 179,949 & 168,470 \\
\quad + MRV / degree       & \phantom{0}82,923 & 114,392 & 117,990 \\
\quad + LCV                & \phantom{0}64,761 & \phantom{0}60,027 & \phantom{0}35,103 \\
\quad + FC + MRV + degree  & \phantom{0}30,066 & \phantom{0}22,951 & \phantom{0}16,588 \\
\bottomrule
\end{tabular}
\\[2pt]
\footnotesize nodes per second on fully censored cells (all runs hit the budget; none solved)
\end{center}

Forward checking's apparent $3.5\times$ node advantage in the all-runs column of
Table~\ref{tab:csp-solved} is therefore not a pruning result at all: it is lower throughput
measured against a shared stopwatch, by $5.4\times$ at $n = 30$, $7.8\times$ at $n = 40$ and
$10.2\times$ at $n = 50$. Under censoring the sign of the metric's meaning
flips --- fewer nodes indicates a \emph{slower} solver, not a better-pruned tree --- and the two
cases are indistinguishable in the reported number. Any scaling curve of nodes against $n$ that
crosses into a censored regime without saying so is reporting throughput and calling it efficiency.

\subsection{Verdict}

Ordering and propagation do what they claim to do, measured in the unit the field uses, on the
instances where that unit is well defined. They also cost $12$--$14\times$ more wall clock than they
save at these sizes, change which instances are solvable in 1 case out of 80 (and that one
adversely), and produce node counts at larger sizes that measure the opposite of what they appear
to. The mechanism that actually changes the cost of solving these instances by orders of magnitude
is the one the family is not named for.


\section{Population-Based Search: The Population Is Not the Mechanism}
\label{sec:pso}

PSO's name attributes its power to a \emph{swarm}: a population of searchers. The canonical
velocity update
\begin{equation}
v \leftarrow \underbrace{w\,v}_{\text{momentum}}
  + \underbrace{c_1 r_1(\text{pbest}-x)}_{\text{cognitive}}
  + \underbrace{c_2 r_2(\text{social}-x)}_{\text{social}},
  \qquad x \leftarrow x + v,
\label{eq:pso}
\end{equation}
contains two candidate mechanisms --- having many searchers, and letting them communicate --- that
the word ``swarm'' conflates. We separate them at a fixed budget of 3{,}030 objective evaluations,
which every arm spends in exactly one instrumented function.

\subsection{The ablation}

Table~\ref{tab:pso-ablation} reports four arms over 30 paired seeds. Arm~(b) sets $c_2 = 0$,
severing communication while retaining thirty searchers; crucially $r_2$ is still drawn, so
arms~(b) and~(c) consume the same random stream and differ only in the mechanism.

\begin{table}[t]
\centering
\small
\caption{Communication ablation on the shuttle-timetabling objective $J$ (minutes of mean student
wait; lower is better). Identical 3{,}030-evaluation budget, 30 paired seeds. Wilcoxon
signed-rank against random search.}
\label{tab:pso-ablation}
\begin{tabular}{@{}lrrl@{}}
\toprule
\textbf{Arm} & \textbf{median $J$} & \textbf{mean $J$ (sd)} & \textbf{vs.\ random search} \\
\midrule
(a) 1 particle, whole budget alone      & 26.10 & 26.61 (2.92) & \textbf{worse}, $p = 7.1\times10^{-5}$ \\
(b) 30 particles, $c_2 = 0$ (no sharing)& 23.79 & 23.57 (1.76) & indistinguishable, $p = 0.22$ \\
(c) 30 particles sharing one \texttt{gbest} & \textbf{15.91} & \textbf{16.23} (1.33) & better, $p < 10^{-8}$ \\
(d) Random search (reference)           & 24.56 & 24.07 (1.61) & --- \\
\bottomrule
\end{tabular}
\end{table}

The answer is unusually clean. Thirty independent searchers are worth \emph{nothing} over a single
random sampler at the same budget: arm~(b) versus arm~(d) gives $p = 0.22$, and the point estimate
favours random search. A single particle spending the entire budget alone is actively
\emph{worse} than random sampling ($p = 7.1\times 10^{-5}$): with no social attractor, inertia
walks it into a corner of the box and it stops sampling usefully. Only when the thirty searchers
share a single best-so-far does the method work at all. Arms~(b) and~(c) differ in one coefficient
and consume the identical random stream, so the paired median difference between them --- $7.61$
minutes, $p = 3.7\times 10^{-9}$ --- is attributable to the social term and to nothing else. In the
units of the problem that is $7.61$ minutes off every student's expected wait, bought by a single
broadcast. Against random search the paired margin is $9.05$ minutes ($p = 1.9\times10^{-9}$).

Neither ``many searchers'' nor ``a searcher with memory'' is the mechanism. The mechanism is a
single shared scalar-plus-vector broadcast once per iteration. This is a quantitative instance of
\citeauthor{sorensen2015metaheuristics}'s objection \citep{sorensen2015metaheuristics}: the
biological metaphor names the population, and the population is the part that does not matter.

\subsection{Three further disconfirmations}

\paragraph{The GA beats the swarm on the swarm's own problem.} At the identical budget, the
real-coded GA reaches median $J = 15.15$ against PSO's $15.91$ (Wilcoxon $W = 106$, $p = 0.008$,
30 paired seeds). We expected the reverse. The diagnosis is in the encoding and is specific: the
objective is invariant under permutation of the departure-time vector, so gene \emph{position}
carries no meaning, while PSO's update~\eqref{eq:pso} moves particle $i$'s coordinate $k$ toward
the attractor's coordinate $k$ --- a correspondence that does not exist. BLX-$\alpha$
\citep{eshelman1993real} has the same exposure but its per-gene interval sampling is less committed
to the alignment, and Gaussian mutation supplies a local-search channel that PSO loses once inertia
has decayed. Both methods nonetheless beat every baseline (Table~\ref{tab:pso-gen}).

\paragraph{The ring topology's celebrated steadiness is not present.} Fully connected
\texttt{gbest} beats the ring on location (mean $J$ $16.23$ vs $17.62$; Wilcoxon $W = 59$,
$p = 1.5\times10^{-4}$). The ring's variance is lower in the direction the literature predicts
\citep{kennedy2002population} ($1.10$ vs $1.77$), but Levene's test does not reject equality of
variance ($p = 0.49$): on 30 seeds the difference is not established, and we decline to claim it.
Both topologies were still improving when the budget expired (last improvement at iteration
$\approx 99$ of 100), so the binding constraint is the evaluation budget, not convergence --- which
also means this comparison says nothing about asymptotic behaviour.

\paragraph{The advantage survives out of sample, but not in the way the objective suggests.}
Schedules optimised on three arrival realisations were scored on 30 held-out realisations
(Table~\ref{tab:pso-gen}). PSO beats the demand-proportional heuristic on \emph{every one} of the
30 held-out days ($p < 10^{-8}$) with an optimism gap of only 2.03 minutes. But the metric a
transport office would recognise --- service level, the fraction of students waiting
$\le 10$ minutes --- is barely improved ($47.1\%$ vs $46.3\%$). What the optimiser actually buys is
a collapse in \emph{stranding}, from $23.0\%$ to $8.8\%$: it carries the evening surge that the
heuristic abandons. A study reporting only service level would have concluded the search was
worthless.

\begin{table}[t]
\centering
\small
\caption{Generalisation to 30 held-out arrival realisations (seeds 100--129). $J$ in minutes.
Service level is the percentage of students waiting $\le 10\,min$.}
\label{tab:pso-gen}
\begin{tabular}{@{}lrrrrr@{}}
\toprule
\textbf{Schedule} & \textbf{train $J$} & \textbf{held-out $J$} & \textbf{optimism gap} &
\textbf{service level} & \textbf{stranded} \\
\midrule
PSO (best of 30 seeds)  & 14.23 & \textbf{16.26} & 2.03  & 47.1\% & \textbf{8.8\%} \\
GA (best of 30 seeds)   & 14.42 & 16.40 & 1.98  & 44.1\% & 7.9\% \\
Demand-proportional     & 23.47 & 22.92 & $-0.54$ & 46.3\% & 23.0\% \\
Uniform headway         & 32.14 & 32.40 & 0.26  & 13.5\% & 9.6\% \\
\bottomrule
\end{tabular}
\end{table}

\paragraph{The constraint that never costs anything.} Across all 27 cells of a
$K \in \{10,\dots,18\} \times B \in \{2,3,4\}$ sweep, the quadratic fleet-overload penalty at the
found optimum is \emph{exactly zero} and peak concurrency never exceeds $B$; at $B = 2$ it
saturates at exactly $2$. The optimiser routes around the constraint rather than paying for it, so
the constraint is invisible in the objective and appears only as a shadow price: cutting the fleet
from three buses to two costs $1.61$ and $1.48$ minutes of mean wait at $K = 10$ and $K = 11$, and
$0.26$ at $K = 12$, while a fourth bus buys nothing at all (differences of $-0.36$ to $-0.66$
minutes, i.e.\ nominally worse, well inside the seed noise of $\mathrm{sd} \approx 1.0$). A third
bus is worth buying; a fourth is not. Reporting penalty terms alone would have reported a
constraint that does not exist.

\section{Reinforcement Learning: The Model Is the Mechanism}
\label{sec:rl}

\subsection{The instance is sound, and every switch is anticipatory}

Value iteration converges in 2,528 sweeps to $\epsilon = 10^{-8}$ with an empirical contraction
rate of $0.98981$ against the theoretical $\gamma = 0.99$ (absolute error $1.93\times10^{-4}$), and
$V^{*}$ agrees with the exact sparse linear solve $(I - \gamma P_\pi)V^{\pi} = R_\pi$ of its own
greedy policy to $4.90\times10^{-9}$. Every policy in this section is scored the same way --- by
exact solution of that linear system under the \emph{true} model --- so rankings carry no
Monte-Carlo noise. Regret is $100\,(V^{*} - V^{\pi})/|V^{*}|$.

The instance has a property worth stating because it makes the lookahead claim provable rather than
merely measured. Delay is charged on the queue \emph{entering} the tick, so it is identical under
both actions; holding discharges the green approach, weakly reducing expected spillback; and
switching pays $c_{\text{switch}} \geq 0$ on top. Therefore
$R(s,\hold) \geq R(s,\swi)$ in every one of the 4,056 states with a genuine choice, with
minimum immediate advantage exactly $+3.000 = c_{\text{switch}}$. Two consequences: the myopic
($\gamma = 0$) policy provably never switches unless maximum-green forces it, so it \emph{coincides
exactly} with always-hold (Table~\ref{tab:rl-baselines} confirms: identical to nine significant
figures); and every voluntary switch in the optimal policy --- 650 states --- is bought
\emph{entirely} out of discounted future value, because the immediate term never once argues for it.

The mirror image locates where the reactive baseline loses: in 1,230 states the optimal policy
holds the green while the \emph{red} queue is longer, and longest-queue-first switches in every one
of them. At $(q_A,q_B) = (11,12)$ with green on $A$ in peak, the reactive rule sees the side road
one vehicle worse off and switches; $V^{*}$ holds, because the main road is mid-discharge and the
clearance tick would throw away five vehicles of capacity to save one vehicle of queue. The optimal
switching curve is emphatically not the diagonal.

\begin{table}[t]
\centering
\small
\caption{Exact regret under the true model. Myopic greedy and always-hold coincide by
construction (see text).}
\label{tab:rl-baselines}
\begin{tabular}{@{}lrrrr@{}}
\toprule
\textbf{Policy} & \textbf{Regret} & \textbf{Action-opt.} & \textbf{Switches/h} & \textbf{Mean delay} \\
\midrule
Value iteration (optimal)      & 0.0\%  & 100.0\% & 84.2 & 26.9\,s/veh \\
Longest-queue-first            & 9.9\%  & 69.5\%  & 97.0 & 29.0\,s/veh \\
Fixed-time (Webster-style)     & 18.1\% & 33.3\%  & 90.0 & 31.7\,s/veh \\
Random (legal actions)         & 43.8\% & 50.1\%  & 97.6 & 38.1\,s/veh \\
Myopic greedy ($\gamma = 0$)   & 45.9\% & 84.0\%  & 51.4 & 40.8\,s/veh \\
Always hold                    & 45.9\% & 84.0\%  & 51.4 & 40.8\,s/veh \\
\bottomrule
\end{tabular}
\end{table}

\subsection{Planning on a learned model beats learning, on identical data}

The decisive comparison holds the data fixed. $Q$-learning generates 400,000 transitions
(4{,}000 episodes of 100 ticks, exploring starts, $\epsilon$-greedy over legal actions, polynomial
step size); the certainty-equivalence planner receives that \emph{literal transition list} --- same
seed, same transitions, same order --- estimates $\hat{P}$ and $\hat{R}$ from it, and plans to
convergence on the estimate. Neither ever touches the true $P$ or $R$; a test greps the learner's
source for model access. State--action coverage is $68.1\%$ of the 11,154 legal pairs, so this
is a claim about a learner that saw two thirds of its problem.

\begin{table}[t]
\centering
\small
\caption{Same 400,000 transitions, same seeds, same order. 5 seeds, $\pm$ standard deviation.}
\label{tab:rl-ce}
\begin{tabular}{@{}lrr@{}}
\toprule
\textbf{Method} & \textbf{Regret} & \textbf{Action-optimality} \\
\midrule
$Q$-learning (default hyperparameters)      & $12.49\% \pm 0.75$ & 65.7\% \\
$Q$-learning (tuned)                        & $2.72\% \pm 0.23$  & --- \\
Certainty-equivalence VI (same data)        & $\mathbf{1.74\%} \pm 0.41$ & \textbf{89.9\%} \\
\bottomrule
\end{tabular}
\end{table}

\paragraph{The tuning step, and why we report it.} The middle row of Table~\ref{tab:rl-ce} is the
most important row in this paper's methodology. A hyperparameter sweep (Appendix~\ref{app:sweep})
shows the \emph{default} $Q$-learning configuration is far from the learner's best: initialising
$Q_0 \approx \bar{V}^{*}$ instead of the nominally optimistic $Q_0 = 0$, together with a step
exponent of $1.0$ instead of $0.7$, moves regret from $12.60\%$ to $4.78\%$ at the sweep's smaller
budget --- the initialisation being the larger share ($12.60\% \to 5.51\%$ on its own). $Q_0 = 0$ is
optimistic here because every reward is negative, so $V^{*} < 0$ everywhere, and that optimism costs
far more than it buys.

Running the comparison against the default learner would have let us report certainty equivalence
beating $Q$-learning $12.49\%$ to $1.74\%$: a sevenfold margin, a cleaner story, and an artifact of
a handicapped baseline. Against the tuned learner the margin is $2.72\%$ versus $1.74\%$ --- modest,
real, and corroborated by the action-optimality column, where the gap is not modest at all
($89.9\%$ versus $65.7\%$). The honest claim is the smaller one, and it is the one we make: on
identical data, planning on the estimated model converts the same transitions into substantially
better decisions than learning the values directly.

\paragraph{Why.} The asymmetry is in what a sample buys. A transition spent on a $Q$-learning
backup updates one state--action pair once, and value propagates backwards only along trajectories
actually walked. The same transition spent on $\hat{P}$ and $\hat{R}$ becomes a permanent part of a
model that unlimited subsequent planning sweeps read from for free --- thousands of sweeps, every
state, no further samples. Once the model is roughly right, extra planning is cheap and extra
experience is not. This is the empirical shadow of the theory relating model estimation to
near-optimal control \citep{kearns2002near, brafman2002rmax, mania2019certainty}.

\subsection{How wrong may the model be?}

If the model is what matters, the operational question is how wrong it may be before learning from
scratch is preferable --- a real question here, because Dhaka does not publish intersection arrival
rates. We planned with assumed arrival rates scaled by $\lambda \in \{0, 0.1, \dots, 2.0\}$ against
the truth and compared to the tuned learner trained on 400,000 real transitions.

The crossing point is stark. Multipliers from $0.1$ to $2.0$ --- rates wrong by a factor of ten too
small or two too large --- all yield regret $\leq 1.79\%$, beating the tuned learner's $2.72\%$. The
\emph{only} model that loses is $\lambda \times 0$: a planner that believes no vehicle ever arrives
($3.58\%$). The crossing therefore sits between multipliers $0.0$ and $0.1$, which is to say:

\begin{quote}
\itshape The model has to be wrong about whether traffic exists at all.
\end{quote}

The reason is structural. The optimal switching curve is driven by the queue dynamics, the discharge
rate, and the clearance cost --- all of which the misspecified planner still has exactly right. The
arrival rates only modulate how far ahead it is worth anticipating. Getting the \emph{structure}
right is what matters; getting the \emph{rates} right barely does. That is a usable thing to be able
to tell a transport office that has no arrival data and never will.

\section{Cross-Family Synthesis}
\label{sec:synthesis}

Four ablations, four misassignments. Taken singly each is a local curiosity. Taken together they
share a structure.

\subsection{The reuse asymmetry}

In every case the component that turned out to be operative is one whose cost is \emph{paid once
and reused many times}, while the eponym's cost is paid per unit of work.

\begin{itemize}[leftmargin=1.4em, itemsep=3pt]
\item \textbf{Search.} The heuristic's information is a global constant obtained by one $O(|E|)$
  reduction. Its cost is per-context; its benefit is per-query. At one query it loses by
  $5.98\times$; past 17 queries it wins. Node counting reports the $n \to \infty$ limit
  unconditionally and never names the assumption.
\item \textbf{Constraints.} Systematic search reconstructs a partial assignment prefix on every
  backtrack; repair carries one complete assignment forward and reuses it across every step. The
  reused object is the assignment itself, which is why min-conflicts pays 0.08\,ms per attempt
  where forward checking pays 1.27\,ms for five fewer nodes.
\item \textbf{Population.} \texttt{gbest} is a broadcast channel: one evaluation, once found,
  informs all thirty particles for the remainder of the run. Sever it and thirty particles decay to
  thirty independent samplers, which is what the $p = 0.22$ against random search measures. The
  population supplies parallelism; only the channel supplies reuse.
\item \textbf{Learning.} A transition spent on a Bellman backup is consumed once. The same
  transition spent on $\hat{P}$ is read by every subsequent planning sweep. Certainty equivalence
  wins because it converts a non-reusable resource into a reusable one.
\end{itemize}

Stated generally: \emph{the mechanisms that survive budget-matched ablation are those that convert
a per-unit cost into an amortised one.} This is not a law, and \S\ref{sec:threats} lists cases
where it would fail. But it predicts the four results here, and it explains the second pattern.

\subsection{Each customary metric declines to charge for reuse}

The metrics are not neutral. Each family's customary metric is denominated in exactly the unit that
makes amortisation invisible (Table~\ref{tab:metrics}). Nodes expanded charges search work but not
heuristic construction. Backtrack counts charge failure but not the cost of the machinery that
avoids it. Iteration or generation counts are denominated per-population, so they hide how the
population is used. Sample counts charge acquisition but not computation on what was acquired ---
which is the whole distinction between $Q$-learning and certainty equivalence.

\begin{table}[t]
\centering
\small
\caption{Each family's customary metric, the resource it does not charge, and the measured
consequence in this paper.}
\label{tab:metrics}
\begin{tabular}{@{}p{2.1cm}p{2.3cm}p{3.4cm}p{5.3cm}@{}}
\toprule
\textbf{Family} & \textbf{Customary metric} & \textbf{What it does not charge} & \textbf{Measured consequence} \\
\midrule
Informed search & nodes expanded & heuristic construction; per-node heuristic evaluation &
  ranking of $\mathrm{A}^*$ vs.\ UCS reverses ($2.15\times$ win $\to$ $5.98\times$ loss);
  benefit overstated $1.48\times$ even excluding setup \\
\addlinespace
Constraint satisfaction & nodes, backtracks & per-node ordering and propagation cost; budget
  censoring &
  orderings save 3--5 nodes and cost $12$--$14\times$ the wall clock; under censoring the metric
  measures throughput, inverting its sign \\
\addlinespace
Population search & iterations, generations & evaluations per iteration; whether the population
  communicates &
  30 non-communicating particles $\equiv$ random search ($p = 0.22$) \\
\addlinespace
Reinforcement learning & samples, episodes & computation performed per sample &
  planning on $\hat{P}$ beats tuned $Q$-learning on identical data ($1.74\%$ vs $2.72\%$) \\
\bottomrule
\end{tabular}
\end{table}

\subsection{A metric-audit rule}

The four cases suggest a checklist that costs little and would have caught all of them:

\begin{enumerate}[leftmargin=1.6em, itemsep=2pt]
\item \textbf{Name the unit.} State what one unit of the reported metric costs in wall clock, and
  verify that it is equal across arms. If a mechanism changes the per-unit cost --- and heuristics,
  propagators, and larger populations all do --- the metric is not comparable across arms without
  reporting the per-unit cost too.
\item \textbf{Charge setup.} Report any per-context precomputation separately, and state the query
  volume at which it amortises. ``$\mathrm{A}^*$ is $2.15\times$ better'' and
  ``$\mathrm{A}^*$ pays off after 17 queries'' are different claims; only the second is actionable.
\item \textbf{Ablate the eponym, not just the hyperparameters.} Sever the named mechanism at fixed
  budget while holding the random stream fixed. In \S\ref{sec:pso} this cost one line
  ($c_2 = 0$, with $r_2$ still drawn) and produced the paper's cleanest result.
\item \textbf{Tune the arm you expect to lose.} \S\ref{sec:rl} would have reported a sevenfold
  effect instead of a $1.6\times$ one had we skipped this. Any ablation that only tunes the
  favoured arm measures tuning effort, not mechanism.
\item \textbf{Check the portfolio for independence.} Before crossing $k$ heuristics with $m$
  algorithms, verify the $k$ heuristics induce distinct orderings (\S\ref{sec:search}). Ours did
  not, and the $3\times3$ grid was a one-dimensional sweep.
\end{enumerate}

\section{Threats to Validity}
\label{sec:threats}

We take these seriously enough to enumerate them at length, because several of our results are
reversals and reversals are fragile.

\paragraph{Implementation constants dominate one result.} The wall-clock reversal in
\S\ref{sec:search} is a statement about \emph{this} implementation: pure Python, NetworkX
adjacency, a cost function evaluated per edge without caching. \citet{burns2012implementing} show
heuristic-search timings varying by an order of magnitude with implementation quality. A compiled
implementation, or one that caches per-edge costs, would shrink the 1.57\,s setup and
move the break-even below 17 queries; a vectorised reduction could plausibly remove it almost
entirely. What is \emph{not} implementation-dependent is the structural point: the setup is
$O(|E|)$ and the search is $O(\text{nodes expanded})$, these scale differently, and node counting
prices the first at zero. The magnitude is ours; the asymmetry is not.

\paragraph{Instances are synthetic.} Every instance is generated, not measured. Parameters are
motivated (class-change cadence, saturation flow of 1,800\,veh/h, 60-minute reneging) but not
calibrated against field data, because the relevant field data for Dhaka is not published --- which
is itself part of the motivation for \S\ref{sec:rl}'s misspecification study. Conclusions about
\emph{mechanisms} should transfer more readily than conclusions about \emph{magnitudes}, but
neither is established for real instances.

\paragraph{One setting, one author, one machine.} All four problems share a domain and an author,
which controls for some confounds and introduces others: a systematic blind spot in problem design
would appear in all four. All timings are from a single laptop-class machine
(Appendix~\ref{app:env}) with no isolation from OS scheduling; timing results should be read as
ratios, not absolutes.

\paragraph{Statistical power.} \S\ref{sec:search} uses 10 query pairs and reports means without
tests --- the effects there are large ($5.98\times$) and the invariance result is exact, but the
timing means carry no error bars. \S\ref{sec:pso} uses 30 paired seeds, adequate for the reported
effects but explicitly \emph{inadequate} for the ring-variance question, which we therefore decline
to claim. \S\ref{sec:rl} uses 5 seeds; the $1.74\%$ vs $2.72\%$ gap is roughly $2.4$ pooled standard
deviations and we do not claim more precision than that.

\paragraph{Scope of the collapse result.} Proposition~\ref{prop:weight} applies to heuristics of the
form $k \cdot \dgeo$. It is not a claim that heuristic portfolios generally collapse; landmark-based,
pattern-database, or learned heuristics are functionally independent by construction. The claim is
that \emph{this} portfolio collapsed, that the collapse was invisible in the study's own reporting,
and that the check is cheap enough to be mandatory.

\paragraph{\S\ref{sec:csp} measures a narrow band of instance sizes, and it is the easy band.}
This is the most serious limitation in the paper and we state it in the strongest form we can. The
24 instances every solver completed have a median search of $13.5$ nodes and finish in under
$1.3$\,ms. On searches that small, any per-node machinery loses on wall clock more or less by
construction: forward checking's fixed cost per assignment cannot be repaid over nine nodes. Our
finding that ordering and propagation cost $12$--$14\times$ what they save is therefore a finding
\emph{about small instances}, and it is exactly the regime where the literature would not claim
propagation helps. The honest scope is: propagation does not pay below $\sim$15 nodes of search,
and our instances do not have a middle regime in which to test where it starts paying, because
they jump from ``solved in 13 nodes'' to ``not solved in five seconds'' with almost nothing
between. Establishing the crossover would need an instance generator with a tunable, gradual
difficulty knob, which ours does not have. What survives regardless of size is the
\emph{censoring-inversion} result (\S5.3), which is a statement about what a node count means when
runs are truncated, not about which solver is better.

\paragraph{Censoring more generally.} At $n \geq 30$ every systematic run exhausts its budget, so
runtimes there are censored and carry no information about the solvers. We confine cross-solver
runtime claims to the uncensored regime and report the large-$n$ columns only to demonstrate the
inversion; they are an anytime comparison, not a scaling law.

\paragraph{Certainty equivalence is favoured by the setting.} \S\ref{sec:rl}'s result holds in a
tabular MDP with 7,098 states, where planning to convergence is free (3.8\,s) and
$\hat{P}$ can be stored exactly. It also depends on a stated modelling choice --- a pessimistic
self-loop for unvisited state--action pairs --- which is favourable to the planner at $68\%$
coverage. Under function approximation, in large or continuous state spaces, or with adversarial
coverage, the comparison could invert. We claim the result for this regime and note that the regime
is the one most classical MDP pedagogy actually occupies.

\paragraph{The synthesis is an interpretation.} \S\ref{sec:synthesis}'s reuse-asymmetry account is
a post-hoc reading of four results, not a tested hypothesis. It generates predictions --- e.g.\ that
a PSO variant broadcasting more information per iteration should outperform \texttt{gbest} by more
than the population size predicts --- which we have not tested. Readers should weight it
accordingly.

\paragraph{No claim of novelty for the individual mechanisms.} That weighted $\mathrm{A}^*$
trades nodes for suboptimality, that min-conflicts is fast, that social learning drives PSO, and
that model-based methods can be sample-efficient are all known. The contribution is the
budget-matched measurement of each on a common testbed, the metric-reversal magnitudes, and the
cross-family structure --- not the discovery of the underlying phenomena.

\section{Conclusion}

We asked whether the mechanism a classical AI method family is named for is the mechanism that
produces its performance, and answered it four times under matched budgets. The heuristic in
heuristic search turned out to be one scalar, priced at zero by the field's own metric, and beaten
outright by an uninformed algorithm. The ordering heuristics in constraint satisfaction reordered
search without changing what was solved. The swarm in particle swarm optimisation turned out to be
a single broadcast channel, with the population itself statistically indistinguishable from random
sampling once the channel was cut. And the learning in reinforcement learning lost, on its own data,
to planning with a model estimated from that data --- and kept losing until the model was wrong
about whether traffic existed.

The common thread is that surviving mechanisms convert per-unit costs into amortised ones, and that
each subfield's customary metric is the one that declines to charge for amortisation. That
coincidence is worth more attention than any individual result here. The metric-audit rule in
\S\ref{sec:synthesis} is our attempt to make the check routine; the released testbed is our attempt
to make it cheap.

\paragraph{Reproducibility.} All code, data-generation scripts, and the scripts that regenerate
every table in this paper are released at
\url{https://github.com/ninadgns/eponym-ablation}. Every experiment is seeded and reproduces bit-for-bit.

\bibliographystyle{plainnat}
\bibliography{refs}

\appendix

\section{Environment}
\label{app:env}

All experiments: Intel Core i5-1135G7 @ 2.40\,GHz (8 threads), 22\,GB
RAM, Linux 7.0.0, Python 3.12.6. Each problem is an independent \texttt{uv} project with a
committed lockfile. Road network: OpenStreetMap extract of the Dhaka bounding box via OSMnx
\citep{boeing2017osmnx}, 55,009 nodes and 137,529 directed edges after simplification.
Test suites: 23 tests (Problem~3) and 28 tests (Problem~4), no skips, both passing at the commit
these results were generated from.

\section{Hyperparameter Sweep for \S\ref{sec:rl}}
\label{app:sweep}

The sweep in Table~\ref{tab:sweep} runs at a smaller budget than the headline comparison in
Table~\ref{tab:rl-ce} --- 1{,}500 episodes (150{,}000 transitions) and 3 seeds rather than
4{,}000 episodes (400{,}000 transitions) and 5 seeds --- so its absolute regrets are higher
throughout and are not directly comparable to Table~\ref{tab:rl-ce}. It is used only to
\emph{select} the tuned configuration, which is then re-run at the full budget.

Three findings are worth stating.

\paragraph{Initialisation, not the step size, is the dominant lever.} $Q_0 = 0$ looks neutral but is
strongly \emph{optimistic} in this MDP: every reward is negative, so $V^{*} < 0$ everywhere and a
zero-initialised table over-values every unvisited action. Setting $Q_0 \approx \bar{V}^{*} = -560$
alone moves regret from $12.60\%$ to $5.51\%$ --- more than any other single change, and more than
the step-size exponent contributes.

\paragraph{The step-size schedule matters, in the direction theory predicts.} At a fixed
150,000-transition budget the polynomial schedule $\alpha_n = (1+n)^{-p}$ substantially outperforms
a constant step: constant $\alpha = 0.1$ and $\alpha = 0.5$ give $24.52\%$ and $21.80\%$ regret
against $12.60\%$ for the baseline schedule. Robbins--Monro conditions say a constant $\alpha$ will
not converge asymptotically; that is \emph{not} what this table measures, since every arm here
stops at the same finite budget. What we observe is worse regret at a fixed budget, which is
consistent with non-convergence but does not demonstrate it.
Within the polynomial family, larger $p$ is better here ($p = 1.0$: $8.72\%$; $p = 0.7$: $12.60\%$;
$p = 0.5$: $17.35\%$).

\paragraph{Exploring starts buy coverage, and coverage buys regret.} Removing exploring starts
collapses state--action coverage from $51.6\%$ to $33.7\%$ and degrades regret to $14.86\%$. Since
coverage bounds what any tabular method can know, we report it beside every learning result rather
than only reporting regret.

The $\gamma$ rows are measured against each $\gamma$'s own $V^{*}$ and so are not comparable to the
others; they are included to show that the learning problem gets harder as the effective horizon
lengthens, which is the expected direction.

\begin{table}[h]
\centering
\small
\caption{$Q$-learning hyperparameter sweep. 15 configurations $\times$ 3 seeds at
150,000 transitions each. Baseline is $\alpha_n = (1+n)^{-0.7}$, $\epsilon$ floor $0.05$,
$Q_0 = 0$, exploring starts on.}
\label{tab:sweep}
\begin{tabular}{@{}lrrr@{}}
\toprule
\textbf{Configuration} & \textbf{Mean regret} & \textbf{sd} & \textbf{Coverage} \\
\midrule
baseline                                   & 12.60\% & 1.46 & 51.6\% \\
\midrule
\multicolumn{4}{@{}l}{\itshape Step-size schedule} \\
$\alpha$ power 0.5                         & 17.35\% & 0.93 & 51.6\% \\
$\alpha$ power 1.0                         &  8.72\% & 0.99 & 51.4\% \\
constant $\alpha = 0.1$                    & 24.52\% & 2.72 & 52.2\% \\
constant $\alpha = 0.5$                    & 21.80\% & 1.04 & 52.0\% \\
\midrule
\multicolumn{4}{@{}l}{\itshape Exploration} \\
$\epsilon$ floor 0.0                       & 12.10\% & 0.86 & 51.5\% \\
$\epsilon$ floor 0.2                       & 11.65\% & 0.22 & 51.4\% \\
$\epsilon$ floor 0.5                       & 10.81\% & 1.00 & 50.6\% \\
no exploring starts                        & 14.86\% & 0.37 & \textbf{33.7\%} \\
\midrule
\multicolumn{4}{@{}l}{\itshape Initialisation} \\
neutral $Q_0 = -560 \approx \bar{V}^{*}$   &  \textbf{5.51\%} & 0.36 & 49.8\% \\
pessimistic $Q_0 = -2000$                  &  9.96\% & 0.21 & 45.6\% \\
\midrule
\textbf{tuned: $Q_0 = -560$, $\alpha$ power 1.0} & \textbf{4.78\%} & 0.14 & 49.8\% \\
\midrule
\multicolumn{4}{@{}l}{\itshape Discount (each measured against its own $V^{*}$; not comparable above)} \\
$\gamma = 0.5$                             &  3.41\% & 0.08 & 52.0\% \\
$\gamma = 0.9$                             &  8.80\% & 0.03 & 51.6\% \\
$\gamma = 0.999$                           & 13.92\% & 1.19 & 51.4\% \\
\bottomrule
\end{tabular}
\end{table}


\end{document}
